\section{本章小结}
先用节点与话题划分职责，再用功能包管理代码和依赖。应熟练完成发布、订阅、定时器和自定义接口的最小闭环。工程运行按“声明依赖、构建安装、加载环境、Launch 启动、观察话题与日志”串联；发生问题时先检查实际使用的包和二进制路径，再检查域、发现范围、接口与 QoS。CLion 使用独立构建目录，运行环境应与终端一致。组件、安全通信和复杂 DDS 配置先会查对应发行版文档，再按项目需要展开。

\Needspace{12\baselineskip}
\begin{problemset}
\renewcommand{\theenumi}{\arabic{enumi}}
\item \label{prob:v2-ros-1}基础：一个包里有两个可执行程序，每个程序创建一个节点。包、进程、节点分别有几个？同一进程能否有多个节点？
\item \label{prob:v2-ros-2}基础：colcon build 后出现 build、install、log。运行前应加载哪个目录中的环境脚本？
\item \label{prob:v2-ros-3}基础：create\_publisher 的整数参数 10 是否表示“每秒发送十次”？哪个机制决定周期发布？
\item \label{prob:v2-ros-4}基础：包名 my\_msgs 与消息文件 My.msg 编译后，C++ 消息类型如何书写？
\item \label{prob:v2-ros-5}提高：发布者和订阅者使用相同话题名，仍收不到数据，按什么顺序排查？
\item \label{prob:v2-ros-6}实践：启动 demo\_nodes\_cpp 的服务端，调用 AddTwoInts，请求 a=-4、b=9，记录响应。
\item \label{prob:v2-ros-7}基础：修改所在发行版的本机发现范围设置后，为什么旧节点可能继续跨机器通信？它与 \nolinkurl{ROS_DOMAIN_ID} 是否能替代认证？
\item \label{prob:v2-ros-8}提高：Launch 中 namespace 为 blue、相对话题 status 重映射为 health，完整话题名是什么？源目录有 YAML，但启动时报文件不存在，应检查什么？
\item \label{prob:v2-ros-9}提高：定时器每 200 毫秒发布消息，INFO 日志按 2000 毫秒限流。两者频率分别是多少？FATAL 是否自动结束程序？
\item \label{prob:v2-ros-10}实践：同一包在 CLion 中运行的是新代码，ros2 run 却仍显示旧文本。说明构建、安装和环境三个环节的排查顺序，并给出修复命令。
\end{problemset}

\section{本章验收}
\begin{enumerate}
    \item 从空目录创建独立工作空间，现场运行字符串发布与订阅；由检查者更改输入文本，展示对应输出。
    \item 在完整工程中增加一个自定义字段并重新构建，展示接口查询、带时间戳的消息和定时器频率。
    \item 使用 rqt 或 Foxglove 展示数据流；关闭发布端再恢复，说明观察到的变化并提交构建命令。
    \item 提交包含清单、构建规则、参数与 Launch 的 training\_tools 包；由检查者指定命名空间和文本，现场启动并展示对应话题。
    \item 现场演示本机通信范围设置与恢复，保存一份独立实验日志；在 IDE 中命中发布回调断点，并说明实际运行的程序路径。
\end{enumerate}
